\section{Moment bounds and convergence of the height field}\label{sec:probability}

The current limits of Section~\ref{sec:chiral} concern insertions at
separated times. We now pass to the height field itself. We first identify its law on smooth neutral tests, then remove endpoint
smoothings to obtain the separated increment moments. The logarithmic
variogram bound extends this identification to finite-energy signed
measures. After these three steps, we prove tightness in negative
regularity, treating ordinary restriction spaces, the outward test
convention, and intrinsic dual spaces separately.

Two estimates drive these passages. An energy-commutator argument gives
uniform ordinary moments for neutral spatial averages. A small neutral
factor separated from all other factors has a correlation that tends to
zero as a positive power of its diameter. The latter estimate makes the
successive removal of endpoint smoothings summable.

Throughout the angular argument, fix \(0<c<2\), and retain the plane state,
\(\kappa\), the vacuum \(\Omega\), the shift \(S=P(i)\), and the
nonnegative endpoint Hamiltonian \(E\) from Section~\ref{sec:angular}.
The Gaussian field appearing before its normalization is identified has
covariance
\begin{equation}\label{eq:prob-log-kernel}
 K_\kappa(x,y)=-2\kappa\log|x-y|.
\end{equation}
We denote this field modulo constants by \(\mathcal H_\kappa\).
All constants below may depend on the fixed parameter and the stated
moment order. None depends on an integer approximation to \(\delta^{-1}\):
\(\delta>0\) is always an arbitrary real mesh size.

\subsection{Ordinary moments from the endpoint Hamiltonian}

On an equal-time column, \(Z_l\) is the diagonal spin at site \(l\), with
values \(\pm1\). These are self-adjoint site operators; they should be
distinguished from marked time rows. A path height increment is the
negative of its step spin. Vacuum moments of the commuting site spins
are their ordinary moments under \(\P_c\).

\begin{lemma}[Ordinary spatial moments]\label{lem:ordinary-moments}
Fix \(0<A<\infty\). Let \(R\ge1\), and let \(g=(g_l)_{l\in\Z}\) be a
finitely supported real sequence such that
\begin{equation}\label{eq:prob-g-conditions}
 \sum_l|g_l|\le AR,\qquad
 \sum_l|g_{l+1}-g_l|\le A,\qquad
 \sum_l|g_{l+1}-g_l|^2\le A/R.
\end{equation}
Then, for every \(1\le b<\infty\),
\begin{equation}\label{eq:prob-ordinary-moments}
 \left\|\sum_lg_lZ_l\right\|_{L^b(\P_c)}\le C_{A,b}.
\end{equation}
In particular the conclusion holds if \(g\) has support of length at
most \(AR\), \(|g_l|\le A\), and
\(|g_{l+1}-g_l|\le A/R\). The same assertions hold with the two axes
interchanged.
\end{lemma}

\begin{proof}
Work throughout this proof in \(\mathcal H_{\mathrm{spin}}\), the closed
span of \(O\Omega\) with \(O\) a finite local spin matrix.
By Lemma~\ref{lem:endpoint-Hamiltonian}, this subspace contains
\(\Omega\) and reduces \(E\). Every finite local spin matrix and its
adjoint preserve the local generating vectors and hence their closed
span. In particular \(D=D_g=\sum_lg_lZ_l\) is bounded and self-adjoint
on \(\mathcal H_{\mathrm{spin}}\). The shift preserves this subspace
in both directions, since it translates local matrices and fixes
\(\Omega\). In this proof \(E\) denotes its self-adjoint restriction
to \(\mathcal H_{\mathrm{spin}}\),
\(B_M=\bbone_{\{E\le M/R\}}\) is its spectral projection there, and
all Hilbert-space norms are taken on this subspace.

First we prove a two-sided compression bound. For fixed \(M\),
Lemma~\ref{lem:endpoint-Hamiltonian} places the range of \(B_M\)
inside the ambient spectral band \(H_0\le C M/R\), when \(R\) is
sufficiently large. The full-space estimates of Lemma~\ref{lem:low-band}
therefore apply to these compressions without increasing their norms.
The identity
\[
 N(i)=\tfrac12S^2(Z_0+Z_1)
\]
and translation by the unitary shift therefore give
\(\|B_M(Z_l+Z_{l+1})B_M\|\le C_M/R\).
Pair the consecutive sites in the sum for \(D\). The coefficients of the
paired sums have total absolute value at most \(AR\); the remaining
single-site terms have total absolute coefficient at most
\(\sum_l|g_{l+1}-g_l|\). Consequently
\begin{equation}\label{eq:prob-two-sided}
 \|B_MDB_M\|\le C_{A,M}.
\end{equation}
For bounded \(R\), this follows instead from
\(\|D\|\le\sum_l|g_l|\).

The local-link description of \(E\) in
Lemma~\ref{lem:endpoint-Hamiltonian} gives
\begin{equation}\label{eq:prob-double-commutator}
 \|[D,[E,D]]\|\le C\sum_l|g_{l+1}-g_l|^2\le C_A/R.
\end{equation}
Indeed the diagonal part of a link commutes with \(D\), and its spin
exchange preserves \(Z_l+Z_{l+1}\); only the coefficient difference
\(g_l-g_{l+1}\) enters each commutator. These identities initially hold
on the local analytic core in \(\mathcal H_{\mathrm{spin}}\), which is
a core for the restricted \(E\). On that core, \([E,D]\) is a finite
sum of bounded local operators. To extend it, approximate a vector in
\(\Dom(E)\subset\mathcal H_{\mathrm{spin}}\) in the graph norm by local
core vectors and use \(EDu=DEu+[E,D]u\) there. Boundedness of \(D\)
and of this commutator, followed by closedness of \(E\), shows that
\(D\) preserves the restricted domain and that the identity extends to
it. The same holds for every fixed local power of \(D\), justifying
the following energy-projection and vacuum calculations on this subspace.

Set \(m_* =\|DB_1\|\), which is finite since \(D\) is bounded.
For a unit vector \(u\in B_1\mathcal H_{\mathrm{spin}}\), expanding the double
commutator yields
\[
 \langle Du,EDu\rangle
 =\tfrac12\langle u,[D,[E,D]]u\rangle
      +\Re\langle Du,DEu\rangle
 \le \frac{C_A+m_*^2}{R}.
\]
Here \(Eu\in B_1\mathcal H_{\mathrm{spin}}\) and \(\|Eu\|\le R^{-1}\).
The component of \(Du\) above energy \(2/R\) thus has squared norm at
most \((C_A+m_*^2)/2\). Its component below \(2/R\) is bounded by
\eqref{eq:prob-two-sided}, with \(M=2\). Taking a supremum gives
\[
 m_*^2\le C_A+\tfrac12m_*^2,
 \qquad\text{and therefore}\qquad
 \|DB_1\|=\|B_1D\|\le C_A.
\]

We next bound powers of \(D\) on the vacuum. Since \(E\Omega=0\),
\[
 2\langle D^k\Omega,ED^k\Omega\rangle
 =\langle\Omega,[D^k,[E,D^k]]\Omega\rangle.
\]
In the site-diagonal representation, let \(D^{\mathrm{sw},l}\) be the
value of \(D\) after exchanging the spins at \(l,l+1\). An exchange
term in this double commutator has, up to a bounded link coefficient
and sign, the multiplier
\(\bigl(D^k-(D^{\mathrm{sw},l})^k\bigr)^2\).
This nonnegative multiplier commutes with the exchange. Furthermore,
\[
 |D-D^{\mathrm{sw},l}|\le2|g_l-g_{l+1}|,
 \qquad
 \bigl|D^k-(D^{\mathrm{sw},l})^k\bigr|^2
 \le C_{A,k}|g_l-g_{l+1}|^2(1+|D|^{2k-2}).
\]
For the second inequality use the factorization of \(a^k-b^k\) and
\(|a-b|\le2A\). Taking the absolute value of an exchange expectation
is bounded by the expectation of its nonnegative multiplier, since the
exchange is a norm-one involution commuting with that multiplier.
Summing the links proves
\begin{equation}\label{eq:prob-power-energy}
 \langle D^k\Omega,ED^k\Omega\rangle
 \le \frac{C_{A,k}}R
       \bigl(1+\|D^{k-1}\Omega\|^2\bigr).
\end{equation}
The low-energy component satisfies
\(\|B_1D^k\Omega\|\le C_A\|D^{k-1}\Omega\|\); on its complement,
positivity of \(E-R^{-1}\) bounds the squared norm by
\(R\langle D^k\Omega,ED^k\Omega\rangle\).
Induction in \(k\) now bounds every \(\|D^k\Omega\|\).
These are the even ordinary moments of \(D\); interpolation gives all
finite orders. Axial symmetry gives the other orientation.
\end{proof}

\begin{corollary}[Neutral blocks]\label{cor:neutral-blocks}
Let \(b\in\R^2\), \(r>0\), and let \(f\) be a measurable real density
supported in a square of side \(Ar\) centered within distance \(Ar\)
of \(b\), with
\[
 \int f=0,\qquad |f|\le A r^{-2}.
\]
For every finite \(q\),
\begin{equation}\label{eq:prob-neutral-block}
 \|\langle h_\delta,f\rangle\|_{L^q}\le C_{A,q}
\end{equation}
uniformly in \(b,r,\delta\). Products of finitely many such blocks have
uniform absolute moment bounds, even when their supports overlap.
\end{corollary}

\begin{proof}
Suppose first \(r/\delta\ge1\), and put \(R=r/\delta\).
Integrate \(f\) on mesh cells, giving masses \(a_{jl}\) of size
\(O(R^{-2})\), with \(O(R)\) possible indices in each direction.
Let \(m_j=\sum_l a_{jl}\), and choose a mass-one discrete profile
\((p_l)\), of width \(O(R)\), bounded by \(C/R\).
The masses \(a_{jl}-m_jp_l\) are neutral in each column. Their discrete
primitives are bounded by \(C/R\), with consecutive differences
\(O(R^{-2})\). Summation by parts and
Lemma~\ref{lem:ordinary-moments}, applied after multiplication by \(R\),
bound the \(L^q\) norm of each column contribution by \(C_q/R\).
Their sum is bounded by Minkowski's inequality.
For the remaining masses \(m_jp_l\), note that
\(\sum_jm_j=0\), \(|m_j|\le C/R\); sum by parts in the other axis,
then average against \(p_l\). This gives the same bound.
If \(r<\delta\), the support meets only a bounded number of neighboring
faces. Neutrality, bounded total variation, and unit height jumps give
a deterministic bound. Finally apply Hölder's inequality to products,
using the required higher moments. No disjointness is used.
\end{proof}

\subsection{Smooth test functions}

\begin{proposition}\label{prop:smooth-tests}
For every finite family of smooth compactly supported densities
\(\phi_1,\ldots,\phi_m\) of integral zero, the vector
\(\bigl(\langle h_\delta,\phi_i\rangle\bigr)_i\) converges jointly in law
and in every joint moment to the corresponding pairings of
\(\mathcal H_\kappa\). The moment convergence is uniform when the tests
vary in a compact smooth family with a common compact support.
\end{proposition}

\begin{proof}
First assume that \(\int\phi(x,y)\,dy=0\) for every \(x\).
For fixed continuum time \(x\), put
\(F(x,y)=\int_{-\infty}^y\phi(x,t)\,dt\).
Discrete summation by parts writes the line pairing as
\[
 \int h_\delta(x,y)\phi(x,y)\,dy
 =-\sum_{l\in\Z}F(x,(l+1)\delta)
       \bigl(h(\lfloor x/\delta\rfloor,l+1)
             -h(\lfloor x/\delta\rfloor,l)\bigr).
\]
The primitive is compactly supported, bounded, and changes by \(O(\delta)\)
at consecutive sites. Lemma~\ref{lem:ordinary-moments} therefore bounds
this pairing in every \(L^q\), uniformly in \(x\) and \(\delta\).
The same remains true for any finite product of line pairings.

Replace the forward height differences in this expression by the
centered differences of Proposition~\ref{prop:gaussian-currents}, shifting
the coefficients accordingly. The coefficient error has size
\(O(\delta)\) and consecutive differences \(O(\delta^2)\), on
\(O(\delta^{-1})\) sites. Dividing this error by \(\delta\) puts it
under Lemma~\ref{lem:ordinary-moments}; the resulting \(L^q\) error is
\(O(\delta)\), uniformly in the line coordinate.

Expand a joint moment as an integral over its time coordinates.
For distinct fixed times, Proposition~\ref{prop:gaussian-currents}
applies to the transverse Riemann sums, locally uniformly on their
compact supports and including the rounding of faces. Their limit is
Wick's formula for transverse derivatives of \(K_\kappa\).
Integration by parts against the compact primitives gives the
undifferentiated test densities and the kernel
\eqref{eq:prob-log-kernel}. The line moment bounds dominate the integrand
uniformly. Equal-time walls have Lebesgue measure zero, so dominated
convergence completes the time integrations. The resulting products of
logarithmic kernels are integrable: in each Wick pairing the two-point
variables are disjoint from those of every other pair, and the logarithm
is locally integrable in two dimensions. Gaussian moments determine the
law. Axial symmetry gives the same result for horizontally line-neutral
tests.

For an arbitrary smooth neutral \(\phi\), let
\[
 q(x)=\int\phi(x,y)\,dy,\qquad
 p_R(y)=R^{-1}p(y/R),\qquad
 \phi_R(x,y)=\phi(x,y)-q(x)p_R(y),
\]
where \(p\) is a fixed smooth compact probability density.
Then \(\phi_R\) is vertically line-neutral, while \(q p_R\) is
horizontally line-neutral. For fixed \(R\), the already proved result
gives
\[
 \lim_{\delta\downarrow0}\E|\langle h_\delta,q p_R\rangle|^2
 =\iint K_\kappa(z,z')q(x)p_R(y)q(x')p_R(y')\,dz\,dz'.
\]
This tends to zero as \(R\to\infty\). To see this, set
\(y=Rt\), \(y'=Rt'\). The additive \(-2\kappa\log R\) contributes
zero because \(\int q=0\). The remaining logarithm is
\(\log\sqrt{((x-x')/R)^2+(t-t')^2}\), which converges almost everywhere
to \(\log|t-t'|\), independently of \(x,x'\).
On the fixed supports it is bounded in absolute value by
\(C+\log_+(1/|t-t'|)\). Dominated convergence and neutrality of \(q\)
prove the assertion. The identical approximation holds for
\(\mathcal H_\kappa\). Thus first taking the mesh limit at fixed \(R\),
and only then taking \(R\to\infty\), proves joint convergence in law.

For any fixed compact support, Corollary~\ref{cor:neutral-blocks} bounds
all higher moments of each fixed smooth pairing. Uniform integrability
therefore upgrades convergence in law to all joint moments. Finally,
small changes in a smooth test on a fixed support give small \(L^q\)
changes by the same corollary, after dividing by the sup norm of the
change. A finite covering of a compact family proves the asserted
uniformity.
\end{proof}

Smooth neutral tests now have the claimed Gaussian limit and all its
joint moments. To recover point increments, we compare each endpoint
with a smooth average at successively smaller scales. The next estimate
controls the correlation contributed by the smallest scale in such a
product.

\subsection{Removing smoothings at separated endpoints}

Let \(\varrho\) be a smooth probability density supported in a small
unit square, and define
\[
 H_{\delta,r}(b)=\int h_\delta(x)r^{-2}
                   \varrho((x-b)/r)\,dx.
\]
Differences \(H_{\delta,r}(b)-H_{\delta,2r}(b)\) obey
Corollary~\ref{cor:neutral-blocks}. At \(r\asymp\delta\), the residual
\(h_\delta(b)-H_{\delta,r}(b)\) is deterministically bounded by unit
height jumps. We will use both kinds of differences as individual
endpoint factors.

\begin{lemma}[A separated endpoint factor]\label{lem:separated-block}
Fix an integer \(k\ge2\), constants \(A,c_0,c_1>0\), a separation
\(d_0>0\), and a bounded region containing all factor centers. Choose
\(r_0=a_kd_0\), with \(a_k>0\) sufficiently small in terms of these
fixed data. Let \(n\le k\) and let
\[
 A_i=\langle h_\delta,\mu_i\rangle,\qquad 1\le i\le n,
\]
where each \(\mu_i\) is a finite signed measure of mass zero, total
variation at most \(A\), and support in a square of side \(Ar_i\)
about \(b_i\). Assume \(\|A_i\|_{L^{2k}}\le A\). These hypotheses
include the neutral density blocks of Corollary~\ref{cor:neutral-blocks}
and the terminal residuals just defined.

Suppose one selected factor has scale \(c_0\delta\le r\le r_0\), all other
scales lie in \([r,r_0]\), and all other centers are at distance at
least \(c_1d_0\) from its center. Then, for some \(\tau>0\),
\begin{equation}\label{eq:prob-separated-block}
 \left|\E\prod_{i=1}^n A_i\right|
 \le C(r/d_0)^\tau.
\end{equation}
The constants are uniform in the centers, the measures, the mesh and
the scales subject to these bounds. The other centers need not be
mutually separated.
\end{lemma}

\begin{proof}
If empty time slabs separate the selected factor from all the others,
the intervening rows damp its short height paths. General endpoint
geometry need not have such a slab in either axial ordering. We will
create one by translating the groups lying to the left and right of
the selected factor's horizontal cluster. Those groups remain separated
from the cluster by horizontal gaps; bounded analytic continuation
across these gaps then returns them, one at a time, to their original
positions.

\smallskip\noindent\textbf{The estimate with empty time slabs.}
First suppose that the selected support has an empty axial slab on each
side, of width at least \(c_2d_0\), separating it from all other factors.
Insert time cuts in these two slabs.
The vectors defined by the exterior factors have bounded Hilbert norms.
Here is the connection with ordinary moments. Read all heights on one
side of a cut relative to a common face, along a common seam and spatial
paths contained on that side. Spatial steps use equal-time diagonal
spins; each time-seam step uses a time spin, with repeated occurrences
allowed when several paths share the step. In a finite homogeneous
system, reflection across the cut transposes the entire product and
complex-conjugates its scalar coefficients. The crossing identity
reverses the time-step spins and preserves the spatial-step spins,
exactly as for the reflected height differences. The squared norm is
therefore the ordinary expectation of the exterior product times its
conjugate reflected product. This finite identity passes to the plane
by Proposition~\ref{prop:physical-plane}. Hölder's inequality and the
ordinary moments of the individual factors bound that expectation.
The factors on a side may overlap; Hölder's inequality requires no
separation among them.

Inside the slab, use neutrality to express the selected factor as an
average of differences from a fixed reference face in its support.
Each difference has a path of \(O(r/\delta)\) steps, with a bounded
number of straight segments. Choose the direction of the difference
so that time runs forward; spatial steps may have either sign.
There remain \(m\ge c_3d_0/\delta\) unmarked time rows on both sides.
Pair consecutive steps along each straight segment. The paired current
identity and Lemma~\ref{lem:low-band} give the sandwiched norm bounds
\begin{equation}\label{eq:prob-damped-costs}
 \begin{aligned}
 \|P(1)^m(\text{two consecutive step sum})P(1)^m\|
       &\le C m^{-1},\\
 \|P(1)^m(\text{one step})P(1)^m\|
       &\le C m^{-\theta_0}
 \end{aligned}
\end{equation}
for some \(0<\theta_0\le1\). To verify the passage from band bounds to
these inequalities, split both sides into the bands
\(H_0\le m^{-1}\) and
\(2^{j-1}/m<H_0\le2^j/m\).
A block between two such bands has current bound a power of the larger
cutoff, while the time rows contribute
\(\exp(-c2^j-c2^{j'})\). The resulting double geometric-exponential
sum is \(O(m^{-1})\) for a paired current. For a single current use the
equal-sign \(O(u)\) and opposite-sign \(O(u^\theta)\) estimates,
with \(\theta_0=\min(1,\theta)\). Bands above a fixed cutoff are
exponentially damped and the uncompressed row norms are bounded.
Unit shifts commute with the damping and do not enlarge these norms.

There are \(O(r/\delta)\) pairs and only a bounded number of unpaired
steps. The bounded total variation of the selected factor and the
exterior norm bounds give
\[
 \left|\E\prod_iA_i\right|
 \le C\left(\frac r{d_0}
               +\left(\frac\delta{d_0}\right)^{\theta_0}\right)
 \le C(r/d_0)^{\tau_0},\qquad \tau_0=\min(1,\theta_0).
\]
This includes a terminal residual: its comparison paths have bounded
length and the same damping margins. Rounding the cuts changes their
widths by only a bounded number of sites.

\smallskip\noindent\textbf{Creating a slab by two translations.}
Cluster the centers
by their first-coordinate projections: consecutive projections belong
to the same cluster if their gap is at most \(a d_0/k\), for a sufficiently
small fixed \(a\). Every cluster has width at most \(a d_0\).
Choose \(a_k\) much smaller than \(a/k\), so that the supports preserve
a gap of size \(c_kd_0\) at each cluster boundary.
In the cluster containing the selected center, every other center has
second-coordinate distance at least \(c_kd_0\) from it. Indeed its
first-coordinate distance is at most \(a d_0\), whereas its Euclidean
distance has the lower bound in the statement.

Translate the clusters strictly to the left as one group vertically,
and do the same independently for the clusters strictly to the right.
For each group there is an interval of translation amounts, of length
at least \(c_kd_0\) and at distance at most \(C_kd_0\) from zero, such
that every translation in the interval leaves a vertical gap of size
\(c_kd_0\) around the selected support. To construct one, exclude the
intervals of translations putting any of its at most \(k\) centers
within a fixed small multiple of \(d_0\) of the selected height.
Their total length is bounded by \(C_kd_0\); choose a larger bounded
translation range and take a complementary interval of the asserted
size. The choices for the two groups are independent. Restrict to
translations by even numbers of lattice sites. Lattice stationarity
preserves each individual \(L^{2k}\) bound under these translations,
so Hölder's inequality still controls the exterior products.
Throughout the resulting rectangle of choices the selected factor has a free slab in the
northward time ordering, so the preceding estimate applies.

\smallskip\noindent\textbf{Returning the two groups.}
The first-coordinate gaps have remained open throughout these vertical
translations. They supply the damping needed to transfer the estimate
back to zero translation. Write
\[
 K=\operatorname{Arg}(S^2),\qquad -\pi\le K\le\pi.
\]
This is bounded spectral functional calculus. The low-band estimate
\(S=\pm1+O(H_0)\) gives \(|K|\le CH_0\) for small \(H_0\); on the
complement the same inequality follows by increasing \(C\), since
\(|K|\le\pi\). Thus
\begin{equation}\label{eq:prob-shift-strip}
 |K|\le CH_0,
 \qquad
 \|P(1)^N e^{izK}\|\le1
       \quad\text{if}\quad |\Im z|\le N/C.
\end{equation}
For integer \(m\), \(e^{imK}=S^{2m}\).
Across a first-coordinate gap containing \(N\ge c_kd_0/\delta\)
unmarked rows, translating an entire exterior group by \(2m\) sites
therefore inserts \(P(1)^NS^{\pm2m}\) between fixed exterior vectors.
This identity follows by conjugating every site and seam insertion
on that side by the site shift, which fixes the vacuum. Reflection and
the ordinary moment estimates bound both exterior vectors uniformly.
Equation~\eqref{eq:prob-shift-strip} extends the pairing to a bounded
holomorphic function of \(z/(d_0/\delta)\) on a strip of fixed positive
width. If one of the exterior groups is absent there is no restoration
step for that group.

For completeness, the elementary analytic estimate used here is the
following. A holomorphic function bounded by one on a fixed-width strip,
whose absolute value is at most \(a\le1\) on a fixed-length real interval
within a fixed bounded distance of zero, has absolute value at zero at
most \(C a^\vartheta\), where \(\vartheta>0\) depends only on those
geometric constants. Take a half-disc based on a shorter part of the
interval. The two-constants theorem bounds the function by
\(a^{\omega}\) on an interior disc, since the harmonic measure of the
base there is bounded below by \(\omega>0\). Cover the remaining path
to zero by finitely many overlapping discs inside the strip.
The three-circles inequality transfers a positive power bound at each
step. The number and radii of these discs depend only on the fixed
strip, interval, and distance bounds.

Our free-slab estimate initially holds only at the allowed integer
values of \(m\). Cauchy's estimate on a slightly smaller strip bounds
the derivative in the normalized real coordinate. The gap between
successive allowed values is \(O(\delta/d_0)\), so throughout a shorter
real interval the bound becomes
\[
 C(r/d_0)^{\tau_0}+C\delta/d_0
 \le C'(r/d_0)^{\tau_0},
\]
using \(r\ge c_0\delta\) and \(\tau_0\le1\).
First restore the left group while keeping the right translation at
any one of its allowed values. The analytic estimate gives the same
bound with a possibly smaller positive exponent, uniformly over those
right translations. Then restore the right group. The two reductions
of the exponent still leave a fixed \(\tau>0\). This proves
\eqref{eq:prob-separated-block} in the original geometry.
\end{proof}

\begin{proposition}\label{prop:offdiagonal}
For every \(k\ge1\), all \(k\)-point height-increment moments converge
to the Wick increment moments of \(\mathcal H_\kappa\), uniformly on
compact sets of tuples whose endpoint sets belonging to different
increments are disjoint.
\end{proposition}

\begin{proof}
The first moment is zero by arrow-reversal symmetry. For \(k\ge2\),
fix such a compact set of tuples and let \(d_0>0\) be a lower bound on
all cross-pair endpoint distances. Endpoints within a single pair may
coincide. Choose \(r_0=a_kd_0\) as in
Lemma~\ref{lem:separated-block}, and put \(r_j=2^{-j}r_0\).
For each endpoint, telescope its bump averages down to the last scale
\(r_J\asymp\delta\), and include the residual from that average to the
specified face value. This is a finite exact identity, including for
points on grid lines under the fixed floor convention.

Expand the product of increments. A factor is either a coarse
increment at scale \(r_0\), or one of the two endpoint blocks at a
dyadic scale, including the terminal residual. All individual blocks
have uniformly bounded moments. The coarse increments also do: join
their endpoints by a bounded-length chain of overlapping scale-\(r_0\)
bumps and telescope, applying Corollary~\ref{cor:neutral-blocks} to
each difference. The number of links is bounded by the fixed compact
size and separation data.

In a summand containing an endpoint block, select one of minimum scale
\(r\). Decompose every other coarse increment by a bump chain avoiding
a disc of radius \(d_0/3\) about the selected endpoint. Such a chain
exists with bounded length because both endpoints of that other pair
are at least \(d_0\) away; a polygonal detour around the disc has length
bounded in terms of the same fixed data. After subdividing into steps
of size \(O(r_0)\), every factor other than the selected one has center
at distance at least \(c d_0\) from it. There are still at most \(k\)
factors in each resulting product; the number of resulting products is
bounded independently of the mesh and dyadic levels. Apply
Lemma~\ref{lem:separated-block} to each of them.

At fixed deepest dyadic index \(j\), there are at most
\(C_k(1+j)^{k-1}\) choices of the other levels and endpoint labels.
Consequently the error from keeping only scales coarser than \(r_{J_0}\)
is bounded by
\begin{equation}\label{eq:prob-desmoothing-tail}
 C\sum_{j\ge J_0}(1+j)^{k-1}2^{-\tau j},
\end{equation}
uniformly once the mesh is fine enough to reach that level. This tends
to zero as \(J_0\to\infty\). Notice that the terminal residual obeys
the same bound; dropping it before applying the estimate would not
justify point evaluation.

For fixed \(J_0\), Proposition~\ref{prop:smooth-tests} gives the joint
moments of the smoothed increments, uniformly in their centers.
The corresponding Gaussian Wick expressions converge uniformly as the
smoothing radius tends to zero: each paired kernel involves endpoints
from two different increments, so its arguments stay a positive
distance apart. Combine this with \eqref{eq:prob-desmoothing-tail}.
If the two endpoints within a pair coincide, that increment and its
smoothed versions are identically zero; no additional exclusion is
needed.
\end{proof}

\subsection{Ordinary finite-energy signed measures}

\begin{lemma}\label{lem:log-integrability}
Let \(\mu\) be a compactly supported finite signed measure for which
\(\iint G(x,y)\,\mu(dx)\mu(dy)\) is an ordinary finite signed-measure
integral. Then \(\mu\) has no atoms, and
\begin{equation}\label{eq:prob-absolute-energy}
 \iint\log_+\frac1{|x-y|}\,|\mu|(dx)|\mu|(dy)<\infty.
\end{equation}
If \(\mu,\nu\) both satisfy this assumption, the corresponding mixed
absolute logarithmic integral is finite as well.
\end{lemma}

\begin{proof}
The total variation of \(\mu\otimes\mu\) is
\(|\mu|\otimes|\mu|\). Ordinary integrability is integrability with
respect to this variation. On a fixed compact support, the negative
part of \(-\log|x-y|\) is bounded, which proves
\eqref{eq:prob-absolute-energy}. An atom would give positive product
mass to the logarithmic diagonal singularity, so there are no atoms.

For the mixed assertion take dyadic squares of side \(2^{-j}\) in a
box containing both supports. Up to an additive bounded kernel,
\(\log_+(1/|x-y|)\) is bounded above by a constant times
\(\sum_{j\ge0}\bbone_{\{|x-y|\le C2^{-j}\}}\).
Write \(a=|\mu|\), \(b=|\nu|\).
At a fixed level the mixed mass of these nearby pairs is at most
\(\sum_Q\sum_{Q'\sim Q}a(Q)b(Q')\), where each square has a bounded
number of neighbors. Cauchy--Schwarz bounds this by
\[
 C\left(\sum_Qa(Q)^2\right)^{1/2}
  \left(\sum_Qb(Q)^2\right)^{1/2}.
\]
The sums of \(\sum_Qa(Q)^2\) and \(\sum_Qb(Q)^2\) over all levels
are finite by their self-integrability: pairs in one square are within
\(\sqrt2\,2^{-j}\), and summing their indicators is bounded by
\(C(1+\log_+(1/|x-y|))\).
A second Cauchy--Schwarz inequality over \(j\) proves the claim.
\end{proof}

\begin{proposition}\label{prop:finite-energy}
For every finite family of compactly supported, neutral, finite signed
measures of the ordinary finite-energy class in
Theorem~\ref{thm:main}, the associated height pairings converge jointly
in law to the centered Gaussian vector with covariance obtained by
integrating \(K_\kappa\) against the pairs of measures.
\end{proposition}

\begin{proof}
Let \(V(a)=\E(h(a)-h(0))^2\) for lattice-face displacement
\(a\in\Z^2\), and put \(V(0)=0\).
Proposition~\ref{prop:variogram} and axial symmetry imply
\[
 V(a)=4\kappa\log(1\vee|a|)+O(1),
\]
with a uniform bounded error over all lattice displacements.
For distinct continuum points \(x,y\) in a fixed compact set, write
\[
 V_\delta(x,y)=
 V\bigl(\lfloor x/\delta\rfloor-\lfloor y/\delta\rfloor\bigr).
\]
For \(0<\delta\le1\) this gives
\begin{equation}\label{eq:prob-rounded-log-bound}
 \left|V_\delta(x,y)-4\kappa\log\delta^{-1}\right|
       \le C\bigl(1+|\log|x-y||\bigr).
\end{equation}
Indeed when \(|x-y|\) exceeds a fixed multiple of \(\delta\), the
lattice separation is comparable to \(|x-y|/\delta\).
In the other case its size is bounded, and
\(\log\delta^{-1}\le C+|\log|x-y||\).
This latter case includes distinct points rounded to the same face.

For any four points, stationarity of increments and zero mean give
\begin{align}\label{eq:prob-four-variograms}
 &\E[(h_\delta(x')-h_\delta(x))
            (h_\delta(y')-h_\delta(y))]\\
 &\qquad=\tfrac12\bigl(
 V_\delta(x',y)+V_\delta(x,y')
       -V_\delta(x',y')-V_\delta(x,y)\bigr).\nonumber
\end{align}
The four copies of \(4\kappa\log\delta^{-1}\) cancel exactly.
Equation~\eqref{eq:prob-rounded-log-bound} therefore supplies a majorant
by a constant and four continuum logarithmic singularities.

For a neutral \(\mu\ne0\), its Jordan parts have the same mass
\(M>0\). The measure
\(M^{-1}\mu^+(dx')\mu^-(dx)\) represents its pairing as an integral
of \(h_\delta(x')-h_\delta(x)\).
Use this representation for two pairings. By
Lemma~\ref{lem:log-integrability}, the four-logarithm majorant is
integrable against the resulting product measure. Geometric coincidences of endpoints belonging to different pairings
have measure zero because the Jordan parts are nonatomic. Distinct
endpoints rounded to a common lattice face remain covered by the same
majorant. Proposition~\ref{prop:offdiagonal} for two increments and dominated
convergence thus give, for every such \(\mu,\nu\),
\begin{equation}\label{eq:prob-energy-covariance}
 \lim_{\delta\downarrow0}
       \E[\langle h_\delta,\mu\rangle
                  \langle h_\delta,\nu\rangle]
     =\iint K_\kappa(x,y)\,\mu(dx)\nu(dy).
\end{equation}
The zero measure causes no exception.

It remains to identify laws beyond second moments. Let
\(\mu_\epsilon=\mu*\varrho_\epsilon\), where the mollifier is smooth,
nonnegative and compactly supported. These measures are smooth and
neutral. Mollification converges in logarithmic energy:
\begin{equation}\label{eq:prob-energy-mollification}
 \iint K_\kappa(x,y)
       (\mu-\mu_\epsilon)(dx)(\mu-\mu_\epsilon)(dy)
       \longrightarrow0.
\end{equation}
To check this for the present ordinary integral, the logarithm convolved
in one or both variables converges pointwise off the diagonal. On the
fixed enlarged compact support its absolute value is bounded by
\(C(1+|\log|x-y||)\), uniformly in small \(\epsilon\).
For \(|x-y|\ge C\epsilon\) this follows by comparison of distances;
for \(|x-y|<C\epsilon\), scaling the bounded compact mollifier gives
\(C+|\log\epsilon|\le C'(1+|\log|x-y||)\).
Lemma~\ref{lem:log-integrability} permits dominated convergence in each
of the four energy terms, proving
\eqref{eq:prob-energy-mollification}.

Apply \eqref{eq:prob-energy-covariance} to
\(\mu-\mu_\epsilon\). It gives
\[
 \lim_{\epsilon\downarrow0}\limsup_{\delta\downarrow0}
 \|\langle h_\delta,\mu-\mu_\epsilon\rangle\|_{L^2}=0.
\]
At fixed \(\epsilon\), Proposition~\ref{prop:smooth-tests} identifies
the Gaussian limit. The same energy estimate approximates the Gaussian
pairing. Applying this argument to a finite family, or to each of its
linear combinations, proves the asserted joint law and cross covariance.
\end{proof}

The scalar pairings and separated increments are now identified. We next
strengthen convergence of individual pairings to convergence in norms.
The ordinary restriction norms use fields localized on a fixed box.
The source's outward seminorm instead tests the global field over an
increasing sequence of boxes; it requires a separate tail estimate.

\subsection{Negative regularity and the outward test convention}

Fix a smooth compactly supported probability density \(\rho\), and write
\(X_\delta=h_\delta-\langle h_\delta,\rho\rangle\).
For a compactly supported smooth cutoff \(\chi\), put
\(X_\delta^\chi=\chi X_\delta\).

\begin{lemma}[Pinned averages]\label{lem:pinned-averages}
Let \(0<\ell\le1\), and let \(\psi_{b,\ell}\) be supported in a square
of side \(A\ell\) about \(b\), with
\(|\psi_{b,\ell}|\le A\ell^{-2}\). For centers \(b\) in a fixed compact
set and every finite \(q\),
\begin{equation}\label{eq:prob-pinned-averages}
 \|\langle X_\delta^\chi,\psi_{b,\ell}\rangle\|_{L^q}
 \le C_{A,q,\chi,\rho}(1+|\log\ell|).
\end{equation}
The estimate is uniform in \(\delta>0\).
\end{lemma}

\begin{proof}
Let \(m=\int\chi\psi_{b,\ell}\), so \(|m|\le C\).
Subtract \(m\) times a smooth mass-one bump of scale \(\ell\) at \(b\).
The difference obeys Corollary~\ref{cor:neutral-blocks}.
Next telescope that bump through the scales
\(\ell,2\ell,\ldots,1\), and finally compare the scale-one bump with
\(\rho\). Each difference is neutral, has bounded total variation,
and is supported and bounded at its stated scale. There are
\(O(1+|\log\ell|)\) differences. The same argument works below the mesh
scale because the final clause of Corollary~\ref{cor:neutral-blocks}
has no lower bound on \(\ell/\delta\).
\end{proof}

\begin{proposition}\label{prop:negative-regularity}
For every bounded open \(U\subset\R^2\) and every \(-1<\alpha<0\),
the pinned fields converge in law, with tightness, in the ordinary
restriction spaces
\[
 \mathcal C^\alpha_{\mathrm{res}}(U):=B^\alpha_{\infty,\infty}(U),\qquad
 B^\alpha_{p,q}(U),\qquad W^{\alpha,p}(U),
 \quad 1\le p<\infty,\quad1\le q\le\infty,
\]
to the corresponding pinning of \(\mathcal H_\kappa\).
Here the Hölder laws are supported on the separable closure of smooth
functions in the displayed norm. Convergence also holds in the literal
outward-scaling completion defined by (16) of
\cite[Definition~2.6(iii)]{DKLM2026}. No boundary regularity of \(U\) is
assumed. An intrinsic dual interpretation of the negative Sobolev norms
is treated at the end of this section.
\end{proposition}

\begin{proof}
Choose a compactly supported \(\chi\) equal to one on a box containing
\(\overline U\) and \(\supp\rho\). Use compactly supported wavelets of
sufficient finite regularity, including their coarse-scale scaling
functions~\cite[Section~4.C]{Daubechies1988}. We use the full-space
Besov coefficient characterization in
\cite[Proposition~2.3]{HairerLabbe2017}. Normalize the scale-\(2^{-j}\) test functions as densities
of size \(O(2^{2j})\), and denote their averages of \(X_\delta^\chi\)
by \(a_{j,k}\). There are \(O(2^{2j})\) relevant indices \(k\).
For \(\alpha<\beta<0\), choose a moment order \(b\) with
\(2+b\beta<0\). Lemma~\ref{lem:pinned-averages} gives
\begin{equation}\label{eq:prob-wavelet-sum}
 \E\sum_{j\ge0}\sum_k|2^{j\beta}a_{j,k}|^b
 \le C_b\sum_{j\ge0}2^{(2+b\beta)j}(1+j)^b<\infty.
\end{equation}
The coarse block has the same bound. This bounds the stronger
\(B^\beta_{\infty,\infty}\) norm in probability. On a fixed compact
support its bounded sets have uniformly vanishing tails in
\(B^\alpha_{\infty,\infty}\), since
\[
 \sup_{j>J,k}2^{j\alpha}|a_{j,k}|
 \le 2^{(\alpha-\beta)J}
          \sup_{j,k}2^{j\beta}|a_{j,k}|.
\]
For finitely many levels only finitely many coefficients remain.
This proves compactness by finite-dimensional truncation, and also puts
the fields and limits in the separable smooth closure.

The same gain proves the other asserted norm compactness. Indeed a
compactly supported distribution bounded in \(B^\beta_{\infty,\infty}\)
has scale blocks satisfying
\(\|\Delta_j X\|_{L^p}\le C2^{-j\beta}\) for every finite \(p\).
Thus
\[
 \sum_{j>J}2^{j\alpha}\|\Delta_jX\|_{L^p}
 \le C\sum_{j>J}2^{j(\alpha-\beta)}\longrightarrow0.
\]
This controls every \(\ell^q\) Besov block norm, including \(q=\infty\).
It also controls the negative Bessel-potential norm: on a scale block,
\((1-\Delta)^{\alpha/2}\) has an \(L^1\)-bounded convolution kernel
after extracting \(2^{j\alpha}\), so summing these bounds works for
all \(p\ge1\). The same summable block bound controls the usual
full-space negative Slobodeckij realization. Passing to its restriction
uses the defining quotient norm. Restricting a global extension to
\(U\) is continuous by definition; no extension operator from \(U\) is
being invoked. Proposition~\ref{prop:smooth-tests} identifies every
subsequential limit, and hence proves the full mesh convergence.

For clarity, retain also the source's printed outward convention.
Let \(Q_0=(-1,1)^2\), and let \(\mathcal T_1(Q_0)\) consist of
compactly supported neutral densities whose zero extension is
one-Lipschitz. For smooth functions its seminorm is
\begin{equation}\label{eq:prob-outward-norm}
 \|f\|_{\mathrm{out},\alpha,U}
 =\sup_{\substack{0<\varepsilon\le1,\ \varphi\in\mathcal T_1(Q_0)\\
                    \supp\varphi\subset U/\varepsilon}}
  \varepsilon^{-\alpha}
     \left|\int f(x/\varepsilon)\varphi(x)\,dx\right|.
\end{equation}
The absolute value is immaterial for the symmetric test family.
Constants vanish in this seminorm. Its completion is understood modulo
its kernel, as in the source.

Write \(\varepsilon=2^{-j}a\), \(1/2\le a\le1\).
The exact mesh identity, including the floor convention, is
\[
 h_\delta(x/\varepsilon)=h_{2^{-j}\delta}(x/a).
\]
After the change of variables \(x=ay\), the test becomes
\(a^2\varphi(ay)\), supported in \(2Q_0\), neutral, with a uniformly
bounded Lipschitz norm. The condition involving \(U\) only reduces this
family. For any fixed \(-1<\gamma<0\), local
\(\mathcal C^\gamma\) control bounds all these pairings. One can see
this directly from a wavelet expansion: cancellation of fine wavelets
against a Lipschitz test gains \(2^{-j}\), and the resulting series
converges because \(1+\gamma>0\).
The preceding moment argument, with a still stronger negative exponent
and arbitrarily high moment order, therefore gives random variables
\(Z_{j,\delta}\) such that
\[
 \sup_\delta\|Z_{j,\delta}\|_{L^b}\le C_b,
 \qquad
 \sup_{2^{-j-1}\le\varepsilon\le2^{-j}}(\text{quantity in }
             \eqref{eq:prob-outward-norm})
       \le C2^{j\alpha}Z_{j,\delta}.
\]
A local pin may be inserted in the rescaled field since every test is
neutral. In particular the supremum over \(j\ge J\) tends to zero in
\(L^b\), uniformly in \(\delta\), by the summability of
\(\sum_{j\ge J}2^{j\alpha}\).
On finitely many bands, the tested original coordinates lie in one
bounded box. Their test norms are bounded by constants depending on the
number of bands, so the stronger local compactness already proved applies.

These estimates also give the required smooth completion, rather than
only a bound on a larger space of functionals. Take a smooth compact
cutoff \(\chi_0\) equal to one on a sufficiently large fixed box and a
compact smooth mollifier \(\theta_\eta\), and use
\[
 f_{\delta,J,\eta}(x)
       =\chi_0(2^{-J}x)(X_\delta*\theta_\eta)(x).
\]
Convolution preserves the local negative Hölder bound, uniformly for
\(0<\eta\le1\), since its rescaled radius on band \(j\) is at most
\(2^{-j}\). For the outward bands \(j\ge J\), the cutoff becomes a bump
of scale at most \(2^{J-j}\) in the rescaled coordinates. The local
\(\mathcal C^\gamma\) test bound gives its contribution the additional
factor \(2^{(J-j)(2+\gamma)}\le1\); comparison of the fixed pin with the
rescaled pin costs at most \(C(1+j)\) by the same dyadic telescoping as
in Lemma~\ref{lem:pinned-averages}. Thus the cutoff approximations have
tails bounded by a constant times
\(\sum_{j\ge J}(1+j)2^{j\alpha}\), uniformly in probability and in the
mollification parameter. On finitely many bands send \(\eta\downarrow0\);
then send \(J\to\infty\). These smooth approximations and the vanishing tail prove
membership and tightness in the completion of
\eqref{eq:prob-outward-norm}. Its limits agree with the smooth-test limit
on each finite band. This proves the literal outward assertion too.
\end{proof}

\subsection{Intrinsic negative Sobolev norms on an arbitrary open set}

The restriction-space proof above already requires no regularity of
\(\partial U\). We give a further argument to cover the convention in
which a negative Sobolev norm is the norm dual to an intrinsic positive
Slobodeckij norm on \(U\). The distinction matters for irregular open
sets, where an extension theorem is unavailable.

For \(0<s<1\), \(1<q<\infty\), write
\[
 \|\phi\|_{W^{s,q}(U)}^q
 =\|\phi\|_{L^q(U)}^q+
   \iint_{U\times U}
       \frac{|\phi(x)-\phi(y)|^q}{|x-y|^{2+sq}}\,dx\,dy.
\]
For \(q=\infty\), use the sup norm and the essential supremum of
\(|\phi(x)-\phi(y)|/|x-y|^s\).
The cells in the following argument are intersections of ordinary
dyadic squares with \(U\); they may have arbitrarily small relative
area and arbitrarily irregular shapes.

\begin{lemma}[Averages on small-area cells]\label{lem:cell-averages}
Let \(Q\) be a square of side \(0<\ell\le1\) in a fixed bounded box,
and let \(A\subset Q\) be measurable with \(|A|>0\). Put
\(t=|A|/\ell^2\). Then for every finite \(b\),
\begin{equation}\label{eq:prob-cell-average}
 \left\|\frac1{|A|}\int_A X_\delta(x)\,dx\right\|_{L^b}
 \le C_b(1+|\log\ell|+|\log t|)
\end{equation}
uniformly in \(\delta\) and in the shape of \(A\).
\end{lemma}

\begin{proof}
We first refine the comparison of line weights. Suppose nonnegative
weights \(a_l\), supported on \(O(R)\) consecutive sites, satisfy
\[
 \sum_l a_l=m>0,\qquad a_l\le M/R,\qquad R\ge1.
\]
Let \(b_l\) be a mass-one profile bounded by \(C/R\), and assume that
the union of the supports of \(a\) and \(b\) lies in one interval of
length \(CR\). We claim
\begin{equation}\label{eq:prob-line-entropy}
 \left\|\sum_l(a_l-mb_l)h_l\right\|_{L^b}
 \le C_bm\bigl(1+\log_+(M/m)\bigr).
\end{equation}
Increasing a fixed constant in \(M\) if needed, assume \(M\ge m\).
For dyadic integer \(v\ge1\), let \(A_va\) replace the weights in
each consecutive block of \(v\) sites by their block average.
The difference \(d=A_va-A_{2v}a\) is neutral in every \(2v\)-block.
For its normalized primitive \(g_l=m^{-1}\sum_{i\le l}d_i\), this gives
\[
 \sum_l|g_l|\le Cv,\qquad
 \sum_l|g_{l+1}-g_l|\le C,\qquad
 \sum_l|g_{l+1}-g_l|^2\le C/v.
\]
The first bound sums a block length times its variation mass.
For the last, use \(\|d\|_\infty\le Cm/v\) and
\(\|d\|_1\le2m\). Lemma~\ref{lem:ordinary-moments} and summation by
parts give cost at most \(C_bm\) for each flattening step. This use of
the lemma explains why its statement retained the three summability
conditions, without a support-diameter hypothesis.

If \(\min(R,Rm/M)>1\), choose a dyadic \(v_0\) comparable from below
to that number; otherwise put \(v_0=1\).
When \(v_0>1\), one can start directly with \(a-A_{v_0}a\).
Its primitive has the same first two bounds, and
\[
 m^{-2}\|a-A_{v_0}a\|_2^2
 \le C m^{-2}\|a\|_2^2
 \le CM/(Rm)\le C/v_0.
\]
When \(v_0=1\), this initial difference vanishes. Continue flattening
until \(v\asymp R\), where comparison with \(mb\) has bounded cost
by the original width-\(R\) case of the lemma. The number of steps is
\(O(1+\log_+(M/m))\), proving \eqref{eq:prob-line-entropy}.

Now rescale the square \(Q\) to a unit square. For any probability
density \(f\le M\) on it, set \(m(x)=\int f(x,y)\,dy\).
Apply \eqref{eq:prob-line-entropy} to its mesh masses on each vertical
line, comparing them with \(m(x)\) times one bounded profile.
Minkowski's inequality bounds the total cost by
\[
 C_b\int m(x)\bigl(1+\log(M/m(x))\bigr)\,dx
 \le C_b(1+\log M).
\]
The last inequality is Jensen's inequality for the concave function
\(u\mapsto u\log(M/u)\), with \(\int m=1\) on an interval of length
one. Compare the remaining horizontal marginal with a bounded profile
by the same line estimate. We have therefore compared the original
density with a fixed scale-\(\ell\) product smoothing, at cost
\(C_b(1+\log M)\). Telescoping that smoothing to scale one and then
to the pin costs \(C_b(1+|\log\ell|)\), as in
Lemma~\ref{lem:pinned-averages}.
If \(\ell<\delta\), all comparisons inside a comparable square use
only a bounded number of faces and have a deterministic bound; the
remaining pin comparison is unchanged. Taking
\(f=\ell^2|A|^{-1}\bbone_A\) in rescaled coordinates gives
\(M=t^{-1}\) and proves \eqref{eq:prob-cell-average}.
\end{proof}

\begin{proposition}[Intrinsic dual convergence]\label{prop:intrinsic-dual}
Let \(U\subset\R^2\) be bounded and open, let \(0<s<1\), and let
\(1\le p<\infty\), with conjugate exponent \(q=p'\).
The pinned fields converge in law, with norm tightness, in
\((W^{s,q}(U))^*\) to the canonically pinned Gaussian field with kernel
\(K_\kappa\). The laws are supported on a separable closed subspace of
this dual. No extension property or regularity of \(\partial U\) is
required.
\end{proposition}

\begin{proof}
Take a fixed square containing \(U\), subdivided dyadically at side
lengths \(\ell_j\asymp2^{-j}\). For a square \(Q\) at level \(j\), put
\[
 U_Q=Q\cap U,\qquad t_Q=|U_Q|/|Q|,\qquad
 X_{\delta,Q}=|U_Q|^{-1}\int_{U_Q}X_\delta.
\]
Ignore cells of zero area. Let
\[
 Z_{\delta,j}=\sum_{Q\text{ at level }j}
                   X_{\delta,Q}\bbone_{U_Q},\qquad
 F_{\delta,j}(\phi)=\int_U Z_{\delta,j}\phi,
\]
and let \(P_j\phi\) denote cell averaging of \(\phi\) on the same
partition. Consistency of averaging implies
\begin{equation}\label{eq:prob-cell-difference}
 (F_{\delta,j+1}-F_{\delta,j})(\phi)
     =\int_U Z_{\delta,j+1}(\phi-P_j\phi).
\end{equation}

Suppose first \(p>1\), so \(q<\infty\), and choose
\(0<\zeta<sq\). Call a parent \(P\) of side \(\ell=\ell_j\) thick
when \(t_P\ge\ell^\zeta\). Jensen's inequality gives
\begin{align*}
 \int_{U_P}|\phi-\phi_P|^q
 &\le |U_P|^{-1}
            \iint_{U_P\times U_P}|\phi(x)-\phi(y)|^q\,dx\,dy\\
 &\le C\ell^{sq}t_P^{-1}
       \iint_{U_P\times U_P}
          \frac{|\phi(x)-\phi(y)|^q}{|x-y|^{2+sq}}\,dx\,dy.
\end{align*}
Summing over thick parents yields
\begin{equation}\label{eq:prob-thick-cell}
 \|\phi-P_j\phi\|_{L^q(\mathrm{thick})}
        \le C\ell^{s-\zeta/q}[\phi]_{W^{s,q}(U)}.
\end{equation}
On all cells, averaging is an \(L^q\) contraction, so the same difference
has \(L^q\) norm at most \(2\|\phi\|_{L^q}\).

Lemma~\ref{lem:cell-averages}, with \(b=p\), implies
\[
 \E\|Z_{\delta,j+1}\|_{L^p(U)}^p\le C_p(1+j)^p.
\]
For this inequality sum
\(|Q|t_Q(1+j+\log(1/t_Q))^p\) over children, and use
\(t(1+\log(1/t))^p\le C_p\) for \(0<t\le1\).
A child of a thin parent has \(t_Q\le4\ell^\zeta\).
The sharper elementary inequality
\[
 t(A+\log(1/t))^p\le C_p A^p t^{1/2},\qquad A\ge1,
\]
therefore gives
\begin{equation}\label{eq:prob-thin-cell}
 \E\|Z_{\delta,j+1}\bbone_{\mathrm{thin}}\|_{L^p(U)}^p
       \le C_p(1+j)^p\ell^{\zeta/2}.
\end{equation}
The sum of the areas \(|Q|\) is bounded by the area of the surrounding
square, independently of the geometry of \(U\).
Equations~\eqref{eq:prob-cell-difference}--\eqref{eq:prob-thin-cell}
and Hölder's inequality now give
\begin{equation}\label{eq:prob-dual-tail-step}
 \E\|F_{\delta,j+1}-F_{\delta,j}\|_{(W^{s,q}(U))^*}
 \le C(1+j)
       \left(\ell_j^{s-\zeta/q}+\ell_j^{\zeta/(2p)}\right).
\end{equation}
Both exponents are positive.

At \(p=1\), \(q=\infty\), the intrinsic Hölder seminorm gives directly
\[
 \|\phi-P_j\phi\|_\infty
       \le C\ell_j^s[\phi]_{C^s(U)}.
\]
There is no need to distinguish thin parents in this case. The first
moment version of the cell estimate gives
\begin{equation}\label{eq:prob-dual-endpoint}
 \E\|F_{\delta,j+1}-F_{\delta,j}\|_{(W^{s,\infty}(U))^*}
       \le C(1+j)\ell_j^s.
\end{equation}
Thus at every finite \(p\) the right-hand side, denoted by \(a_j\),
is summable over \(j\), uniformly in \(\delta\).

Let \(\mathcal E_0\) be the closed linear span in
\(\mathcal E=(W^{s,q}(U))^*\) of the countable collection of cell
indicator pairing functionals. It is a separable Banach space.
For each fixed mesh, \(X_\delta\) is a bounded function on the bounded
set \(U\), and its conditional cell averages converge to it in
\(L^p(U)\). Since \(L^p(U)\) maps continuously into \(\mathcal E\),
they also converge in \(\mathcal E_0\). Summing the preceding
estimates proves
\begin{equation}\label{eq:prob-dual-uniform-tail}
 \sup_\delta\E\|X_\delta-F_{\delta,J}\|_{\mathcal E}
       \le\sum_{j\ge J}a_j\longrightarrow0.
\end{equation}

Every coefficient \(X_{\delta,Q}\) is the pairing with the neutral
measure
\[
 \mu_Q=|U_Q|^{-1}\bbone_{U_Q}\,dx-\rho\,dx.
\]
This bounded compact density has ordinary finite logarithmic energy,
regardless of the shape or area of \(U_Q\). Proposition~\ref{prop:finite-energy}
therefore gives joint convergence of the coefficients at each finite
level. Couple their Gaussian limits using the countable family of
Gaussian pairings with \(\mu_Q\), and denote the corresponding
finite-level elements of \(\mathcal E_0\) by \(F_J\).
Finite-dimensional convergence and the lower semicontinuity of the norm
transfer the bound
\(\E\|F_{j+1}-F_j\|_{\mathcal E}\le a_j\).
The sum of these norms is almost surely finite and integrable, so
\(F_J\) converges almost surely and in \(L^1(\mathcal E_0)\) to
an \(\mathcal E_0\)-valued random variable \(F\).
Finite-level convergence and \eqref{eq:prob-dual-uniform-tail}
then prove \(X_\delta\Rightarrow F\) in \(\mathcal E_0\), including
norm tightness.

We finally identify the limit on the whole intrinsic test space, not
only on interior smooth tests. For \(\phi\in W^{s,q}(U)\), the measure
\[
 \mu_\phi=\phi\bbone_U\,dx-\left(\int_U\phi\right)\rho\,dx
\]
has ordinary finite logarithmic energy. Indeed \(q>1\) or \(q=\infty\),
and on a fixed compact set Hölder's inequality for the logarithmic
kernel gives
\[
 \iint |\log|x-y||\,|f(x)f(y)|\,dx\,dy
       \le C_q\|f\|_{L^1}\|f\|_{L^q}.
\]
For finite \(q\), \(P_J\phi\to\phi\) in \(L^q(U)\); for \(q=\infty\)
the intrinsic Hölder bound gives uniform essential convergence.
The associated neutral measures consequently converge in logarithmic
energy. Hence \(F_J(\phi)\) converges to the canonical pinned Gaussian
pairing with \(\mu_\phi\). This identifies \(F\) on every intrinsic
test and rules out an unidentified boundary functional in the larger
dual. The proof is complete.
\end{proof}

\begin{remark}[Other difference-norm conventions]
If a negative Besov convention is dual to a positive intrinsic Besov
difference norm, choose \(0<t<s<1\). The positive test norm of order
\(s\) controls the intrinsic \(W^{t,q}\) norm: in a dyadic difference
expansion, the gain \(2^{-j(s-t)}\) is summable, for every sequence
exponent including its endpoints. The convergence just proved for
\((W^{t,q}(U))^*\) therefore gives the corresponding negative Besov
dual convergence of order \(-s\). This is in addition to the ordinary
restriction Besov conclusion of Proposition~\ref{prop:negative-regularity}.
\end{remark}

For \(0<c<2\), the section has established all three convergence mechanisms with
covariance \(K_\kappa\). Since
\(K_\kappa=4\pi\kappa\,G\), the squared multiplier is
\(4\pi\kappa\). Proposition~\ref{prop:normalization} computes it for
\(c<1\); the conditional normalization theorem quoted in the
introduction applies after these convergence statements have been
proved for \(1\le c<2\).

\subsection{The endpoint topology implication}

The endpoint argument uses a different route to Gaussian convergence.
We isolate the additional implication needed for the ordinary restriction
spaces and arbitrary-domain intrinsic duals, rather than identify these
spaces with the outward seminorm in \cite{DKLM2026}.

\begin{proposition}[Endpoint topology implication]\label{prop:endpoint-topologies}
At \(c=2\), suppose that the height pairings with every finite family
of ordinary finite-energy neutral measures converge jointly to those
of \(\sigma\Gamma\). Then the commonly pinned fields also converge,
with tightness, in all the negative-regularity spaces stated in
Propositions~\ref{prop:negative-regularity} and~\ref{prop:intrinsic-dual}.
This includes the ordinary restriction H\"older, Besov and Sobolev
spaces on every bounded open set, the outward-scaling completion,
and the explicitly defined intrinsic Sobolev duals.
\end{proposition}

\begin{proof}
We obtain the required moment bounds directly from
\cite[Theorem~4.5 and equations~(28)--(30)]{DKLM2026}, which applies
at \(c=2\). For \(2n\) increments whose continuum endpoints lie in
a fixed bounded box, that theorem implies, for \(0<\delta\le1\),
\begin{equation}\label{eq:prob-endpoint-log-moments}
 \left|\E_2\prod_{i=1}^{2n}
       \bigl(h_\delta(x_i)-h_\delta(y_i)\bigr)\right|
 \le C_n\sum_\pi\prod_{\{i,j\}\in\pi}
 \left(1+\sum_{u\in\{x_i,y_i\}}
              \sum_{v\in\{x_j,y_j\}}
                    \log_+\frac1{|u-v|}\right).
\end{equation}
Here and below inequalities with a logarithmic singularity are used
off the corresponding coincidence diagonals. To check the mesh cutoff,
let \(d_{ij}\) be the minimum cross-endpoint distance after rounding
to lattice faces, and \(D_{ij}\) the continuum distance. Rounding gives
\(D_{ij}\le\delta(d_{ij}+C)\), while each lattice increment length
is at most \(C/\delta\). The logarithmic branch of the cited bound
therefore satisfies
\[
 1+\log_+\frac{C/\delta}{1\vee d_{ij}}
       \le C'\left(1+\log_+\frac1{D_{ij}}\right).
\]
Its other branch is at most one. The logarithm of the minimum of
four distances is bounded by the sum of their positive logarithms,
giving \eqref{eq:prob-endpoint-log-moments}. If a rounded increment
is zero, the entire moment vanishes; otherwise every scale ratio is
well-defined. Thus the estimate also
handles distinct continuum endpoints that round to the same face.

Let \(f\) be a probability density supported in a square of side
\(0<\ell\le1\) and bounded by \(C/(\ell^2t)\), where \(0<t\le1\).
For every point \(v\), the disc-mass bound gives
\begin{align*}
 \int\log_+\frac1{|u-v|}\,f(u)\,du
 &=\int_0^\infty\!\left(\int_{|u-v|<e^{-s}}f(u)\,du\right)ds\\
 &\le\int_0^\infty\min\left(1,
                      \frac{C e^{-2s}}{\ell^2t}\right)ds
 \le C'\bigl(1+|\log\ell|+|\log t|\bigr).
\end{align*}
Write \(\langle X_\delta,f\rangle\) as the integral of
\(h_\delta(x)-h_\delta(y)\) against \(f(x)\rho(y)\,dx\,dy\).
After expanding its \(2n\)-th moment, each pair in
\eqref{eq:prob-endpoint-log-moments} involves endpoint variables
disjoint from those in every other pair. The pair integrals therefore
multiply. The fixed smooth pin has a uniform logarithmic integral,
and the last display bounds the other ones. Consequently
\[
 \E_2|\langle X_\delta,f\rangle|^{2n}
 \le C_n\bigl(1+|\log\ell|+|\log t|\bigr)^n.
\]
This is stronger than the bound needed in
\eqref{eq:prob-cell-average}. Taking
\(f=|A|^{-1}\bbone_A\), for an arbitrary positive-area subset
\(A\) of the square, proves that estimate at the endpoint.
Signed densities bounded by \(C\ell^{-2}\) are dominated by a
constant times the uniform probability density on their supporting
square. The same moment expansion therefore proves
\eqref{eq:prob-pinned-averages}, including its smooth cutoff.
Interpolation gives all finite moment orders. Coarser meshes have
uniform deterministic bounds on the fixed box.

We also need scale-uniform neutral blocks for the outward completion.
If both endpoints of every increment lie in a square of side \(r\),
the preceding rounding argument replaces
\(1+\log_+(1/D_{ij})\) by
\(1+\log_+(r/D_{ij})\). Integration against densities bounded by
\(C r^{-2}\) gives a constant independent of \(r\) and \(\delta\)
when \(\delta\le r\). When \(\delta>r\), unit jumps give the
same bound deterministically. For a neutral signed density \(g\)
on that square, choose a probability density \(p\le C r^{-2}\)
there and use the exact representation
\[
 \langle h_\delta,g\rangle
   =\iint\bigl(h_\delta(x)-h_\delta(y)\bigr)g(x)p(y)\,dx\,dy.
\]
Domination by \(|g(x)|p(y)\le C r^{-4}\) gives the required
moment bound without dividing by a possibly small Jordan mass.
This proves the neutral-block estimate
\eqref{eq:prob-neutral-block} at \(c=2\).

We may now repeat the wavelet proof of
Proposition~\ref{prop:negative-regularity}: for
\(\alpha<\beta<0\) choose \(b\) with \(2+b\beta<0\).
The endpoint pinned-average bound gives
\eqref{eq:prob-wavelet-sum}; finite-level truncations and the gain
\(2^{j(\alpha-\beta)}\) yield compactness in every asserted
restriction norm. Smooth neutral pairings identify the limits by
the hypothesis. The global localization and continuous restriction
to \(U\) require no boundary extension theorem. The neutral-block
bound and the exact mesh-rescaling identity give the outward tails
and smooth approximations by the same proof.

Finally the endpoint cell-average estimate gives every thick- and
thin-cell bound in the proof of Proposition~\ref{prop:intrinsic-dual},
including the case \(p=1\), \(q=\infty\). Its finite-level
coefficients are pairings with ordinary finite-energy measures, so
their joint convergence is part of the hypothesis. The summable dual
norm tail proves convergence in the same separable closed subspace;
energy approximation identifies the limit on the entire intrinsic
test space. This proves the proposition.
\end{proof}
